\chapter{Topologie des espaces métriques}\label{ch:b2:metric}

La topologie de la droite réelle (volume de première année) se
généralise, presque sans changer un mot, à tout ensemble muni d'une
distance. Le gain est énorme~: suites de fonctions, matrices, courbes
--- toutes deviennent des points d'\omterm{def:b2:metric:def}{espaces métriques}, et les trois
piliers démontrés ici --- la complétude avec le théorème du point fixe
de Banach, la compacité, la connexité --- s'appliquent à elles
uniformément. Ce chapitre est la colonne vertébrale de toute la moitié
analytique du livre.

\section{Espaces métriques}

\begin{definition}\label{def:b2:metric:def}
Un \emph{espace métrique}\index{espace métrique} est un ensemble $X$
muni d'une application $d \colon X \times X \to \R_+$ telle que, pour
tous $x, y, z$~:
\[
d(x,y) = 0 \iff x = y,
\qquad
d(x,y) = d(y,x),
\qquad
d(x,z) \leq d(x,y) + d(y,z).
\]
Boules~: $B(a, r) = \{x : d(a,x) < r\}$ (ouverte), $\overline B(a,r) =
\{x : d(a,x) \leq r\}$ (fermée). Une partie $A \subseteq X$ devient un
espace métrique pour la distance induite.
\end{definition}

\begin{example}\label{ex:b2:metric:examples}
$\R$ avec $\abs{x - y}$~; $\R^n$ avec l'une des distances
\[
d_1(x,y) = \sum_i \abs{x_i - y_i},
\quad
d_2(x,y) = \Bigl(\sum_i (x_i - y_i)^2\Bigr)^{1/2},
\quad
d_\infty(x,y) = \max_i \abs{x_i - y_i};
\]
l'ensemble $C(\intcc{a}{b})$ des fonctions \omterm{def:b2:metric:continuity}{continues} avec la
\emph{distance sup} $d_\infty(f, g) = \sup_{\intcc{a}{b}} \abs{f -
g}$ (finie~: $f - g$ est bornée)~; tout ensemble muni de la distance
\emph{discrète} ($d(x,y) = 1$ pour $x \neq y$). Les distances issues
de normes font l'objet du~\cref{ch:b2:nvs}.
\end{example}

\begin{definition}[Topologie d'un espace métrique]\label{def:b2:metric:topology}
$U \subseteq X$ est \emph{ouvert}\index{ouvert!dans un espace métrique}
lorsque chaque point de $U$ est le centre d'une boule contenue dans
$U$~; $F$ est \emph{fermé} lorsque son complémentaire est ouvert. Les
voisinages, l'intérieur, l'adhérence, la densité, la frontière se
définissent exactement comme sur la droite réelle (volume de première
année), les boules remplaçant les intervalles, et les énoncés qui y
sont démontrés --- réunions/intersections d'ouverts, caractérisations
de l'intérieur et de l'adhérence, l'adhérence comme plus petit fermé
contenant --- se transposent avec les mêmes démonstrations. Les boules
ouvertes sont ouvertes, les boules fermées sont fermées (inégalité
triangulaire).
\end{definition}

\begin{example}[Intérieur, adhérence, frontière sur un ensemble]\label{ex:b2:metric:intclosure}
Dans $\R$, soit $A = \intoc{0}{1} \cup \{2\}$. Intérieur~:
$\intoo{0}{1}$ --- autour de tout $x \in \intoo01$ une petite boule
reste dans $A$~; autour de $1$, toute boule $\intoo{1-r}{1+r}$ fuit
de $A$ vers la droite, donc $1$ n'est pas intérieur~; et le point
isolé $2$ ne l'est pas non plus. Adhérence~:
$\intcc{0}{1} \cup \{2\}$ (le point $0$ est une limite de $A$,
rien d'autre n'est ajouté). Frontière (adhérence privée de
l'intérieur)~: $\{0, 1, 2\}$. Notez les asymétries à retenir~: une
extrémité peut appartenir à un ensemble sans être intérieure ($1$),
peut être adhérente sans y appartenir ($0$), et un point isolé est
sa propre frontière ($2$). La même comptabilité s'applique mot pour
mot dans tout \omterm{def:b2:metric:def}{espace métrique}, avec des boules à la place des
intervalles.
\end{example}

\begin{definition}[Limites, continuité]\label{def:b2:metric:continuity}
$x_n \to x$ dans $X$ lorsque $d(x_n, x) \to 0$. Une application $f
\colon X \to Y$ entre \omterm{def:b2:metric:def}{espaces métriques} est
\emph{continue}\index{continu} en $a$ lorsque
\[
\forall \varepsilon > 0,\ \exists\delta > 0,\quad
d_X(x, a) \leq \delta \implies d_Y\bigl(f(x), f(a)\bigr) \leq
\varepsilon ;
\]
de façon équivalente (même démonstration que sur $\R$), $f(x_n) \to
f(a)$ pour toute suite $x_n \to a$. $f$ est
\emph{lipschitzienne}\index{lipschitzien} de constante $k$ lorsque
$d_Y(f(x), f(y)) \leq k\, d_X(x, y)$ toujours --- alors uniformément
continue, donc continue.
\end{definition}

\begin{theorem}[Caractérisation globale de la continuité]\label{thm:b2:metric:globalcontinuity}
$f \colon X \to Y$ est \omterm{def:b2:metric:continuity}{continue} (en tout point) si et seulement si
l'image réciproque de tout \omterm{def:b2:metric:topology}{ouvert} est ouverte --- si et seulement si
l'image réciproque de tout fermé est fermée.
\end{theorem}

\begin{proof}
($\Rightarrow$) Soit $V \subseteq Y$ \omterm{def:b2:metric:topology}{ouvert} et $a \in f^{-1}(V)$~:
une certaine boule $B(f(a), \varepsilon) \subseteq V$~; la continuité
en $a$ fournit $\delta$ tel que $f(B(a, \delta)) \subseteq B(f(a),
\varepsilon)$, donc $B(a, \delta) \subseteq f^{-1}(V)$.

($\Leftarrow$) Étant donné $a$ et $\varepsilon$~: $f^{-1}\bigl(B(f(a),
\varepsilon)\bigr)$ est \omterm{def:b2:metric:topology}{ouvert} et contient $a$, donc contient une
boule $B(a, \delta)$~: c'est la définition de la continuité en $a$.
Fermés~: complémentaires (\cref{prop:b2:structures:images}).
\end{proof}

\section{Espaces complets}

\begin{definition}\label{def:b2:metric:complete}
Une suite $(x_n)$ est \emph{de Cauchy} lorsque $\sup_{p, q \geq N}
d(x_p, x_q) \to 0$ quand $N \to \infty$. Un \omterm{def:b2:metric:def}{espace métrique} est
\emph{complet}\index{espace métrique complet} lorsque toute suite de
Cauchy converge. Convergente $\Rightarrow$ de Cauchy toujours~; les
parties fermées d'espaces complets sont complètes, et les parties
complètes d'un espace quelconque sont fermées (mêmes démonstrations
que sur $\R$~: volume de première année).
\end{definition}

\begin{example}[Cauchy sans limite]\label{ex:b2:metric:cauchynolimit}
Dans $X = \Q$ avec la distance usuelle, les troncatures décimales de
$\sqrt2$,
\[
x_0 = 1,\quad x_1 = 1.4,\quad x_2 = 1.41,\quad x_3 = 1.414,
\quad\dots
\]
vérifient $\abs{x_p - x_q} \leq 10^{-\min(p,q)}$~: de Cauchy dans
$\Q$. Une limite dans $\Q$ serait aussi la limite dans $\R$, à
savoir $\sqrt2 \notin \Q$~: aucune limite n'existe dans $X$.
L'incomplétude est la présence de telles «~limites fantômes~»~; la
complétude de $\R$ a été construite dans le volume de première année
précisément pour donner un foyer à toute suite de Cauchy.
\end{example}

\begin{theorem}\label{thm:b2:metric:rncomplete}
$\R^n$ (l'une quelconque des trois distances du~%
\cref{ex:b2:metric:examples}) et $\bigl(C(\intcc{a}{b}),
d_\infty\bigr)$ sont \omterm{def:b2:metric:complete}{complets}.
\end{theorem}

\begin{proof}
$\R^n$~: une suite de Cauchy est de Cauchy dans chaque coordonnée
(chaque $\abs{x_i - y_i} \leq d(x,y)$ pour les trois distances), donc
chaque coordonnée converge (complétude de $\R$, volume de première
année), et la convergence coordonnée par coordonnée entraîne la
convergence pour $d_\infty$ (nombre fini de coordonnées), donc pour
les trois (les trois distances se dominent mutuellement à des
facteurs constants près~: $d_\infty \leq d_2 \leq d_1 \leq
n\,d_\infty$).

$C(\intcc{a}{b})$~: soit $(f_n)$ de Cauchy pour $d_\infty$. Pour
chaque $x$, $(f_n(x))$ est de Cauchy dans $\R$ ($\abs{f_p(x) -
f_q(x)} \leq d_\infty(f_p, f_q)$)~: converge vers un $f(x)$. En
passant à la limite dans $\abs{f_p(x) - f_q(x)} \leq \varepsilon$
(valable pour $p, q \geq N_\varepsilon$, tout $x$) quand $q \to
\infty$~: $\abs{f_p(x) - f(x)} \leq \varepsilon$ pour tout $x$,
c'est-à-dire $d_\infty(f_p, f) \leq \varepsilon$~: convergence
uniforme. La limite est \omterm{def:b2:metric:continuity}{continue}~: étant donné $\varepsilon$,
choisir $p$ avec $\sup\abs{f_p - f} \leq \varepsilon$, puis utiliser
la continuité de $f_p$ en $a$ et le découpage en trois termes
\[
\abs{f(x) - f(a)} \leq \abs{f(x) - f_p(x)} + \abs{f_p(x) - f_p(a)}
+ \abs{f_p(a) - f(a)} \leq 3\varepsilon
\]
pour $x$ proche de $a$. (Cet «~argument des $3\varepsilon$~» revient
comme théorème de la limite uniforme du~\cref{ch:b2:funcseq}.)
\end{proof}

\begin{example}[Ouverts et fermés reconnus par la continuité]\label{ex:b2:metric:recognize}
La caractérisation globale
(\cref{thm:b2:metric:globalcontinuity}) est l'outil quotidien de la
comptabilité topologique. Dans $\R^2$~: l'ensemble $\{(x, y) : x^2 +
y^2 < 1,\ y > x^3\}$ est \omterm{def:b2:metric:topology}{ouvert} --- c'est $g^{-1}(\intoo{-\infty}{1})
\cap h^{-1}(\intoo{0}{+\infty})$ pour les fonctions \omterm{def:b2:metric:continuity}{continues} $g(x,y)
= x^2 + y^2$ et $h(x, y) = y - x^3$, une intersection de deux images
réciproques d'\omterm{def:b2:metric:topology}{ouverts}. Dans $\bigl(C(\intcc01), d_\infty\bigr)$~:
l'ensemble des fonctions avec $f(0) = f(1)$ et $\int_0^1 f = 0$ est
fermé --- l'image réciproque de $\{(0,0)\}$ par l'application \omterm{def:b2:metric:continuity}{continue}
$f \mapsto \bigl(f(0) - f(1),\ \int_0^1 f\bigr)$ à valeurs dans $\R^2$
(chaque coordonnée est $1$-lipschitzienne, comme dans le~%
\cref{exo:b2:metric:3}). La méthode ne dessine jamais~: exhiber une
application \omterm{def:b2:metric:continuity}{continue}, lire l'ensemble comme une image réciproque,
invoquer le théorème.
\end{example}

\begin{example}[Un fermé défini par une infinité de conditions]\label{ex:b2:metric:lipclosed}
Dans $\bigl(C(\intcc01), d_\infty\bigr)$, l'ensemble
\[
L = \{f : \abs{f(x) - f(y)} \leq \abs{x - y}
\ \text{pour tous } x, y\}
\]
des fonctions $1$-lipschitziennes est fermé, bien qu'il soit
découpé par une infinité non \omterm{def:b2:structures:countable}{dénombrable} de conditions~: pour chaque
couple fixé $(x, y)$, l'application $f \mapsto \abs{f(x) - f(y)} -
\abs{x - y}$ est \omterm{def:b2:metric:continuity}{continue} (les évaluations sont
$1$-lipschitziennes), donc chaque condition prise seule définit un
fermé, et $L$ est l'\emph{intersection} de cette famille --- une
intersection quelconque de fermés est fermée. Le même modèle certifie
la fermeture pour les fonctions monotones, les fonctions convexes,
les fonctions majorées par une $g$ fixée~: les limites uniformes
héritent de toute propriété exprimable comme une famille de
contraintes ponctuelles fermées. Ce dont les limites uniformes
n'héritent \emph{pas} automatiquement --- la dérivabilité, par
exemple --- est exactement ce pour quoi le~\cref{ch:b2:funcseq} devra
peiner.
\end{example}

\begin{theorem}[Théorème du point fixe de Banach]\label{thm:b2:metric:banach}
Soit $X$ un \omterm{def:b2:metric:complete}{espace métrique complet} non vide et $f \colon X \to X$
une \emph{contraction}~: \omterm{def:b2:metric:continuity}{lipschitzienne} de constante $k < 1$. Alors
$f$ a un unique point fixe $\ell$, et toute orbite $x_{n+1} = f(x_n)$
converge vers $\ell$, avec\index{théorème du point fixe de Banach}
\[
d(x_n, \ell) \leq \frac{k^n}{1 - k}\, d(x_1, x_0) .
\]
\end{theorem}

\begin{proof}
Unicité~: deux points fixes sont à distance $\leq k$ fois
elle-même. Existence~: $d(x_{n+1}, x_n) \leq k^n d(x_1, x_0)$ par
récurrence, donc pour $q > p$,
\[
d(x_q, x_p) \leq \sum_{j=p}^{q-1} d(x_{j+1}, x_j)
\leq d(x_1, x_0) \sum_{j \geq p} k^j
= \frac{k^p}{1-k}\, d(x_1, x_0) \xrightarrow[p\to\infty]{} 0 :
\]
de Cauchy, donc convergente vers un $\ell$~; la continuité de $f$
passe $x_{n+1} = f(x_n)$ à la limite~: $\ell = f(\ell)$. La majoration
de l'erreur est l'estimation affichée avec $q \to \infty$.
\end{proof}

\begin{example}[Une équation intégrale]\label{ex:b2:metric:integralequation}
Sur $X = C(\intcc{0}{1})$ (\omterm{def:b2:metric:complete}{complet},
\cref{thm:b2:metric:rncomplete}), considérons $T(f)(x) = 1 + \frac12
\int_0^x f(t)\,\dd t$. Pour $f, g \in X$~:
\[
\abs{T(f)(x) - T(g)(x)} \leq \frac12 \int_0^x \abs{f - g} \leq
\frac12\, d_\infty(f, g),
\]
donc $T$ est une $\frac12$-contraction~: elle a un unique point fixe
\omterm{def:b2:metric:continuity}{continu} --- la solution de $f' = \frac f2$, $f(0) = 1$, à savoir
$\eu^{x/2}$. Ce schéma, industrialisé, devient le théorème de
Cauchy--Lipschitz du~\cref{ch:b2:diffeq}.
\end{example}

\begin{example}[Un point fixe numérique~: $x = \cos x$]\label{ex:b2:metric:cosfix}
Sur le \omterm{def:b2:metric:complete}{complet} $X = \intcc{0}{1}$, l'application $f = \cos$ envoie
$X$ dans $\intcc{\cos 1}{1} \subseteq X$ et est une contraction~: par
l'inégalité des accroissements finis,
\[
\abs{\cos x - \cos y} \leq \bigl(\sup_{\intcc01}\abs{\sin}\bigr)
\abs{x - y} = (\sin 1)\abs{x - y},
\qquad \sin 1 \approx 0.841 < 1 .
\]
Banach~: une unique solution de $x = \cos x$ dans $\intcc01$ (donc
dans $\R$~: tout point fixe réel est dans $\intcc{-1}{1}$, puis dans
$\intcc{\cos 1}{1}$ après une application), et l'itération $x_{n+1} =
\cos x_n$ converge vers elle depuis tout point de départ~: $x_\infty
\approx 0.739085$, le fameux nombre obtenu en martelant la touche
cosinus d'une calculatrice. La majoration de l'erreur prédit une
décroissance en $(\sin1)^n/(1 - \sin1)$ --- environ un chiffre tous
les $13$ appuis~; la majoration a posteriori du problème de fin de
semaine de ce chapitre (question 14) certifie chaque étape à la
volée.
\end{example}

\section{Compacité}

\begin{definition}\label{def:b2:metric:compact}
Un \omterm{def:b2:metric:def}{espace métrique} $X$ est \emph{compact}\index{espace métrique
compact} lorsque toute suite de $X$ admet une sous-suite convergeant
\emph{dans} $X$ (la propriété de Bolzano--Weierstrass). Une partie est
compacte lorsqu'elle l'est pour la distance induite.
\end{definition}

\begin{theorem}[Premières propriétés]\label{thm:b2:metric:compactprops}
\begin{enumerate}
  \item Une partie compacte est fermée et bornée~; une partie fermée
        d'un espace \omterm{def:b2:metric:compact}{compact} est compacte.
  \item Dans $\R^n$, la réciproque est vraie~: \omterm{def:b2:metric:compact}{compact} $\iff$ fermé
        et borné.
  \item Une image \omterm{def:b2:metric:continuity}{continue} d'un espace \omterm{def:b2:metric:compact}{compact} est compacte~; une
        fonction réelle \omterm{def:b2:metric:continuity}{continue} sur un espace \omterm{def:b2:metric:compact}{compact} non vide est
        bornée et atteint ses bornes.
  \item (Heine) Une application \omterm{def:b2:metric:continuity}{continue} sur un espace \omterm{def:b2:metric:compact}{compact} est
        uniformément \omterm{def:b2:metric:continuity}{continue}.
  \item Produits~: si $X, Y$ sont \omterm{def:b2:metric:compact}{compacts}, alors $X \times Y$ l'est
        (avec $d\bigl((x,y),(x',y')\bigr) = d(x,x') + d(y,y')$).
\end{enumerate}
\end{theorem}

\begin{proof}
(1) Mêmes arguments que sur la droite (volume de première année)~:
une suite s'échappant à l'infini ou convergeant à l'extérieur n'a
aucune sous-suite convergeant à l'intérieur~; pour la seconde
assertion, extraire dans le \omterm{def:b2:metric:compact}{compact} ambiant et utiliser la fermeture.

(2) Les suites bornées de $\R^n$ ont des sous-suites convergentes
composante par composante~: extraire sur la première coordonnée
(Bolzano--Weierstrass sur $\R$), puis, de cette sous-suite, sur la
deuxième, et ainsi de suite ($n$ extractions successives)~; la
fermeture garde la limite à l'intérieur.

(3) Étant donné $(f(x_n))$, extraire $x_{\varphi(n)} \to x \in X$~;
la continuité donne $f(x_{\varphi(n)}) \to f(x) \in f(X)$. Cas réel~:
la compacité de $f(X) \subseteq \R$ la rend fermée et bornée, et
$\sup f(X) \in f(X)$ (le sup d'un ensemble lui est adhérent, et
$f(X)$ est fermé).

(4) La démonstration de première année se transfère mot pour mot~;
la voici, en habit métrique. Supposons $f \colon X \to Y$ \omterm{def:b2:metric:continuity}{continue}
sur le \omterm{def:b2:metric:compact}{compact} $X$ mais non uniformément \omterm{def:b2:metric:continuity}{continue}~: un certain
$\varepsilon > 0$ admet, pour tout $n$, des points avec
\[
d_X(x_n, y_n) \leq \frac{1}{n+1}
\qquad\text{et}\qquad
d_Y\bigl(f(x_n), f(y_n)\bigr) > \varepsilon .
\]
Extraire $x_{\varphi(n)} \to a \in X$~; alors $y_{\varphi(n)} \to a$
aussi (les distances mutuelles tendent vers $0$). La continuité en
$a$ envoie les deux suites images vers $f(a)$, donc
$d_Y\bigl(f(x_{\varphi(n)}), f(y_{\varphi(n)})\bigr) \to 0$ ---
contredisant l'écart uniforme $> \varepsilon$. La compacité a fourni
exactement une chose~: le point d'accumulation $a$ où appliquer la
continuité ordinaire.

(5) Extraire sur les coordonnées de $X$, puis à nouveau sur les
coordonnées de $Y$.
\end{proof}

\begin{example}[Le théorème de Heine, avec et sans compacité]\label{ex:b2:metric:heinework}
Sur $\intcc{0}{1}$, la fonction $x \mapsto x^2$ est uniformément
\omterm{def:b2:metric:continuity}{continue} --- Heine le dit sans aucun calcul, mais l'estimation
directe est instructive~:
\[
\abs{x^2 - y^2} = \abs{x + y}\,\abs{x - y} \leq 2\abs{x - y},
\]
donc $\delta = \varepsilon/2$ convient \emph{pour tous les points à
la fois}. Sur $\R$ la même fonction n'est pas uniformément \omterm{def:b2:metric:continuity}{continue}~:
avec $x_n = n$ et $y_n = n + \frac1n$, l'écart $\abs{x_n - y_n} =
\frac1n \to 0$ tandis que $\abs{x_n^2 - y_n^2} = 2 + \frac1{n^2} \geq
2$~: aucun $\delta$ unique ne sert $\varepsilon = 1$. Le mécanisme
est visible~: la constante de Lipschitz locale $\abs{x + y}$ est
bornée sur un \omterm{def:b2:metric:compact}{compact} et non bornée sur $\R$ --- le théorème de Heine
est exactement l'énoncé que la compacité plafonne uniformément de
telles constantes locales.
\end{example}

\begin{method}[Prouver qu'un ensemble est compact]\label{met:b2:metric:compactness}
Trois voies, par \omterm{def:b2:structures:generated}{ordre} de fréquence. (1) \emph{Reconnaissance
ambiante~:} dans $\R^n$ (ou tout espace normé de dimension finie,
\cref{ch:b2:nvs}), vérifier fermé --- typiquement comme image
réciproque, \cref{ex:b2:metric:recognize} --- et borné. (2)
\emph{Hérédité~:} une partie fermée d'un \omterm{def:b2:metric:compact}{compact} connu est compacte~;
une réunion finie ou un produit de \omterm{def:b2:metric:compact}{compacts} est \omterm{def:b2:metric:compact}{compact}~; une image
\omterm{def:b2:metric:continuity}{continue} d'un \omterm{def:b2:metric:compact}{compact} est compacte. (3) \emph{À mains nues~:}
extraire une sous-suite convergente d'une suite quelconque ---
généralement par extractions successives coordonnée par coordonnée.
Pour prouver la \emph{non}-compacité, un seul témoin suffit~: une
suite sans sous-suite convergente, le plus souvent des points à
distance mutuelle $\geq \varepsilon$.
\end{method}

\begin{example}[Distances entre ensembles~: la compacité gagne son pain]\label{ex:b2:metric:setdistance}
Soit $K$ \omterm{def:b2:metric:compact}{compact}, $F$ fermé, $K \cap F = \emptyset$ dans un \omterm{def:b2:metric:def}{espace
métrique}. Alors
\[
d(K, F) = \inf\,\{d(x, y) : x \in K,\ y \in F\} > 0 :
\]
la fonction $x \mapsto d(x, F)$ est \omterm{def:b2:metric:continuity}{continue}
(\cref{exo:b2:metric:11}) et positive sur $K$ ($d(x, F) = 0$
mettrait $x \in \overline F = F$), donc elle atteint un minimum
positif sur le \omterm{def:b2:metric:compact}{compact} $K$
(\cref{thm:b2:metric:compactprops}~(3)). La compacité n'est pas
décorative~: pour deux ensembles \emph{fermés} l'infimum peut
s'annuler sans être atteint --- dans $\R^2$, l'hyperbole $F_1 = \{xy
= 1\}$ et l'axe $F_2 = \{y = 0\}$ sont deux fermés disjoints avec
$d(F_1, F_2) = 0$ (les points $(n, \frac1n)$ approchent l'axe). La
fuite à l'infini est exactement ce que la compacité interdit.
\end{example}

\begin{theorem}[Borel--Lebesgue]\label{thm:b2:metric:borellebesgue}
Un \omterm{def:b2:metric:def}{espace métrique} $X$ est \omterm{def:b2:metric:compact}{compact} si et seulement si tout
recouvrement de $X$ par des \omterm{def:b2:metric:topology}{ouverts} admet un sous-recouvrement
\emph{fini}.\index{théorème de Borel--Lebesgue}
\end{theorem}

\begin{proof}
($\Leftarrow$) Supposons que $(x_n)$ n'a aucune sous-suite
convergente. Nous affirmons que tout $x \in X$ a une boule $B(x,
r_x)$ contenant $x_n$ pour un nombre fini d'indices $n$ seulement~:
sinon, toute boule $B(x, \frac1{k+1})$ contiendrait une infinité de
termes, et en choisissant des indices
\[
\varphi(0) < \varphi(1) < \varphi(2) < \cdots
\quad\text{avec}\quad
x_{\varphi(k)} \in B\Bigl(x, \frac{1}{k+1}\Bigr)
\]
(possible à chaque étape précisément parce qu'il reste une infinité
de candidats) on construirait une sous-suite convergeant vers $x$.
Les boules $B(x, r_x)$ recouvrent $X$~; si un nombre fini d'entre
elles recouvrait $X$, l'ensemble d'indices $\N$ serait une réunion
finie d'ensembles finis~: absurde.

($\Rightarrow$) Deux étapes. \emph{Nombre de Lebesgue~:} pour un
recouvrement \omterm{def:b2:metric:topology}{ouvert} $(U_i)$ d'un \omterm{def:b2:metric:compact}{compact} $X$, il existe $\rho > 0$
tel que toute boule de rayon $\rho$ est contenue dans un $U_i$.
Sinon, pour chaque $n$ choisir $x_n$ avec $B(x_n, \frac{1}{n+1})$
dans aucun $U_i$~; extraire $x_{\varphi(n)} \to x \in U_{i_0}
\supseteq B(x, r)$~; pour $n$ grand, $B(x_{\varphi(n)},
\frac{1}{\varphi(n)+1}) \subseteq B(x, r) \subseteq U_{i_0}$~:
contradiction. \emph{Précompacité~:} pour tout $\rho > 0$, un nombre
fini de boules de rayon $\rho$ recouvre $X$. Sinon choisir par
récurrence $x_{n+1}$ hors de $B(x_0, \rho) \cup \dots \cup B(x_n,
\rho)$~: la suite a des distances deux à deux $\geq \rho$, donc
aucune sous-suite de Cauchy --- donc aucune convergente~:
contradiction. En combinant~: recouvrir $X$ par un nombre fini de
boules de rayon $\rho$ (le nombre de Lebesgue), chacune dans un
certain $U_i$~: un sous-recouvrement fini.
\end{proof}

\begin{example}[Un $\varepsilon$-réseau, compté]\label{ex:b2:metric:epsnet}
La précompacité (issue de la démonstration du~%
\cref{thm:b2:metric:borellebesgue}) est très concrète sur
$\intcc{0}{1}$~: pour $\varepsilon > 0$, les $\lceil
\frac{1}{2\varepsilon}\rceil$ boules centrées en $\varepsilon,
3\varepsilon, 5\varepsilon, \dots$ de rayon $\varepsilon$ le
recouvrent --- environ $\frac1{2\varepsilon}$ boules, et aucun
recouvrement ne peut se contenter de moins de $\frac{1}{2\varepsilon}$
d'entre elles (chaque boule couvre une longueur au plus
$2\varepsilon$). Dans $\intcc01^2$ le compte s'élève au carré, à
l'\omterm{def:b2:structures:generated}{ordre} $\varepsilon^{-2}$~: les nombres de recouvrement croissent
comme $\varepsilon^{-d}$ en dimension $d$ --- une face quantitative
de la compacité, et la raison pour laquelle les boules unités de
dimension infinie du~\cref{ch:b2:nvs} (où aucun
$\frac13$-réseau fini n'existe) ne peuvent être \omterm{def:b2:metric:compact}{compactes}.
\end{example}

\begin{example}[Lire la compacité sur les recouvrements]\label{ex:b2:metric:covers}
L'intervalle \omterm{def:b2:metric:topology}{semi-ouvert} $\intoc{0}{1}$ est recouvert par les \omterm{def:b2:metric:topology}{ouverts}
$U_n = \intoo{\frac1n}{2}$, $n \geq 1$~; toute sous-famille finie a
un plus grand indice $N$ et manque $\intoc{0}{\frac1N}$~: aucun
sous-recouvrement fini, donc $\intoc{0}{1}$ n'est pas \omterm{def:b2:metric:compact}{compact} --- ce
que la définition séquentielle voit à travers $x_n = \frac1n$, dont
la limite $0$ s'échappe. En revanche, ajouter le seul point $0$
répare les deux diagnostics d'un coup~: sur $\intcc{0}{1}$ tout tel
recouvrement doit contenir un ensemble contenant $0$, qui avale un
segment initial entier, et un nombre fini d'ensembles finit le reste.
Les deux langages du~\cref{thm:b2:metric:borellebesgue} échouent ou
réussissent toujours de concert --- les recouvrements détectent la
fuite exactement là où les suites la détectent.
\end{example}

\begin{remark}[Perspectives dans ce volume]\label{rem:b2:metric:perspectives}
Ce chapitre est le mur porteur du volume~; observez où chaque pilier
porte du poids. \emph{Complétude}~: le critère de Cauchy devient le
test de convergence des séries dans les espaces de Banach
(\cref{ch:b2:series}), la convergence uniforme du~\cref{ch:b2:funcseq}
est exactement la convergence dans le \omterm{def:b2:metric:complete}{complet} $\bigl(C,
d_\infty\bigr)$, et Cauchy--Lipschitz (\cref{ch:b2:diffeq}) est le
théorème du point fixe de Banach déguisé en équation intégrale.
\emph{Compacité}~: elle prouve l'équivalence des normes
(\cref{ch:b2:nvs}), l'atteinte des extrema pour l'optimisation du~%
\cref{ch:b2:diffcalc}, et l'existence des meilleures approximations
(le problème de fin de semaine du~\cref{ch:b2:nvs}). \emph{Connexité}~:
elle globalise les énoncés locaux --- l'unicité des solutions
d'équations différentielles, le théorème des valeurs intermédiaires
sur les courbes (\cref{ch:b2:curves}), et les deux composantes de
$GL_n(\R)$ que la théorie de l'orientation (\cref{ch:b2:multint})
gardera séparées.
\end{remark}

\begin{remark}[Pièges classiques]\label{rem:b2:metric:pitfalls}
(i) «~Fermé et borné implique \omterm{def:b2:metric:compact}{compact}~» est un théorème sur $\R^n$,
non sur les \omterm{def:b2:metric:def}{espaces métriques}~: un ensemble infini muni de la
distance discrète est fermé et borné en lui-même mais non \omterm{def:b2:metric:compact}{compact}
(\cref{exo:b2:metric:4}), et la boule unité fermée de $C(\intcc01)$
échoue elle aussi (\cref{ch:b2:nvs}). (ii) La complétude est une
propriété de la \emph{distance}, non de la topologie~: $\R$ avec
$d(x,y) = \abs{\arctan x - \arctan y}$ a les mêmes suites convergentes
usuelles mais est incomplet (\cref{exo:b2:metric:1}). (iii) Une
bijection \omterm{def:b2:metric:continuity}{continue} n'est pas nécessairement un homéomorphisme --- la
paramétrisation du cercle du~\cref{exo:b2:metric:7}~; la compacité de
la source y remédie. (iv) Le théorème de Banach a besoin de $k < 1$
\emph{uniformément}~: la condition $d(f(x), f(y)) < d(x,y)$ à elle
seule ne garantit rien sur un espace non \omterm{def:b2:metric:compact}{compact}
(\cref{exo:b2:metric:5}). (v) \omterm{def:b2:metric:connected}{Connexe} n'implique pas \omterm{def:b2:metric:connected}{connexe} par arcs
en général --- mais pour les parties \omterm{def:b2:metric:topology}{ouvertes} des espaces normés
rencontrées dans ce livre, les deux coïncident (\cref{ch:b2:nvs}).
\end{remark}

\begin{remark}[Où ce chapitre est utilisé]
Partout dans la moitié analytique. La complétude de
$C(\intcc{a}{b})$ alimente les théorèmes de convergence du~%
\cref{ch:b2:funcseq} et la théorie de Cauchy--Lipschitz du~%
\cref{ch:b2:diffeq} (le problème de fin de semaine de ce chapitre
prouve déjà le théorème local de Picard--Lindelöf)~; la compacité
donne l'équivalence des normes en dimension finie (\cref{ch:b2:nvs})
et l'existence des extrema dans le~\cref{ch:b2:diffcalc}~; la
connexité sous-tend les arguments de valeurs intermédiaires du~%
\cref{ch:b2:realfun} et la globalisation de l'unicité pour les
équations différentielles. Dans le volume de troisième année, la
compacité dans les espaces de fonctions (le théorème
d'Arzelà--Ascoli) et le théorème de Baire (\cref{exo:b2:metric:12}
ici) deviennent des outils d'usage quotidien.
\end{remark}

\begin{omfigure}
\begin{tikzpicture}[xscale=6.4, yscale=0.8]
  \draw[omDef, line width=3pt] (0,3) -- (1,3);
  \draw[omDef, line width=3pt] (0,2) -- (0.3333,2);
  \draw[omDef, line width=3pt] (0.6667,2) -- (1,2);
  \foreach \a in {0, 0.6667} {
    \draw[omDef, line width=3pt] (\a,1) -- (\a + 0.1111,1);
    \draw[omDef, line width=3pt]
      (\a + 0.2222,1) -- (\a + 0.3333,1);
  }
  \foreach \a in {0, 0.2222, 0.6667, 0.8889} {
    \draw[omThm, line width=3pt] (\a,0) -- (\a + 0.037,0);
    \draw[omThm, line width=3pt]
      (\a + 0.0741,0) -- (\a + 0.1111,0);
  }
  \node[left, font=\small] at (0,3) {$C_0$};
  \node[left, font=\small] at (0,2) {$C_1$};
  \node[left, font=\small] at (0,1) {$C_2$};
  \node[left, font=\small] at (0,0) {$C_3$};
\end{tikzpicture}

{\small Les premières étapes de l'ensemble de Cantor
(\cref{exo:b2:metric:8})~: chaque niveau supprime le tiers central
\omterm{def:b2:metric:topology}{ouvert} de chaque segment. L'intersection $C = \bigcap_n C_n$ est
compacte, d'intérieur vide et de longueur nulle, et pourtant
équipotente à $\R$ --- et elle revient comme \emph{point fixe} d'une
contraction sur les ensembles dans le problème de fin de semaine de
ce chapitre (question 22).}
\end{omfigure}

\section{Connexité}

\begin{definition}\label{def:b2:metric:connected}
$X$ est \emph{connexe}\index{espace connexe} lorsqu'il n'admet aucune
partition en deux parties \omterm{def:b2:metric:topology}{ouvertes} non vides --- de façon
équivalente, lorsque ses seules parties à la fois \omterm{def:b2:metric:topology}{ouvertes} et fermées
sont $\emptyset$ et $X$. $X$ est \emph{connexe par arcs} lorsque deux
points quelconques sont joints par une application \omterm{def:b2:metric:continuity}{continue} $\gamma
\colon \intcc{0}{1} \to X$.
\end{definition}

\begin{theorem}\label{thm:b2:metric:connectedness}
\begin{enumerate}
  \item Les parties \omterm{def:b2:metric:connected}{connexes} de $\R$ sont exactement les intervalles.
  \item Une image \omterm{def:b2:metric:continuity}{continue} d'un \omterm{def:b2:metric:connected}{espace connexe} est \omterm{def:b2:metric:connected}{connexe} --- d'où
        le théorème des valeurs intermédiaires général~: une fonction
        réelle \omterm{def:b2:metric:continuity}{continue} sur un \omterm{def:b2:metric:connected}{espace connexe} prend toute valeur
        comprise entre deux de ses valeurs.
  \item \omterm{def:b2:metric:connected}{Connexe} par arcs $\Rightarrow$ \omterm{def:b2:metric:connected}{connexe}. (La réciproque est
        fausse en général~; elle est vraie pour les parties \omterm{def:b2:metric:topology}{ouvertes}
        des espaces normés, \cref{ch:b2:nvs}.)
\end{enumerate}
\end{theorem}

\begin{proof}
(1) Un $A$ non intervalle manque un $z$ entre deux de ses points~:
$A = (A \cap \intoo{-\infty}{z}) \cup (A \cap \intoo{z}{+\infty})$ le
scinde en deux morceaux non vides \omterm{def:b2:metric:topology}{ouverts} (dans $A$). Réciproquement,
soit $I$ un intervalle et $I = U \cup V$ une partition en parties non
vides relativement \omterm{def:b2:metric:topology}{ouvertes}~; choisir $a \in U$, $b \in V$, disons $a
< b$, et poser $s = \sup\,(U \cap \intcc{a}{b})$, un point de
$\intcc{a}{b} \subseteq I$. Si $s \in U$~: alors $s \neq b$, et
l'ouverture relative de $U$ met tout un intervalle autour de $s$
(intersecté avec $I$) dans $U$ --- donc des points de $U \cap
\intcc{a}{b}$ dépassent $s$, contredisant le supremum. Si $s \in V$~:
l'ouverture relative de $V$ met un intervalle $\intoo{s - r}{s + r}
\cap I$ dans $V$~; mais le supremum est adhérent à $U \cap
\intcc{a}{b}$, qui doit rencontrer cet intervalle --- contradiction
avec $U \cap V = \emptyset$. (C'est l'argument des ouverts-fermés de
première année pour $\R$, mené à l'intérieur de $I$.)

(2) Si $f(X) = U' \cup V'$ se scinde en parties non vides relativement
\omterm{def:b2:metric:topology}{ouvertes}, alors $X = f^{-1}(U') \cup f^{-1}(V')$ scinde $X$
(\cref{thm:b2:metric:globalcontinuity}). TVI~: $f(X) \subseteq \R$
est \omterm{def:b2:metric:connected}{connexe}, donc un intervalle par (1).

(3) Supposons $X = U \cup V$, tous deux non vides \omterm{def:b2:metric:topology}{ouverts}, et joignons
$a \in U$ à $b \in V$ par un chemin $\gamma$~: alors $\gamma^{-1}(U),
\gamma^{-1}(V)$ scindent $\intcc{0}{1}$, contredisant (1).
\end{proof}

\begin{example}\label{ex:b2:metric:glnr}
$GL_n(\R)$ n'est pas \omterm{def:b2:metric:connected}{connexe}~: $\det$ est \omterm{def:b2:metric:continuity}{continu} (un polynôme en les
coefficients) sur $\R^*$, qui n'est pas \omterm{def:b2:metric:connected}{connexe}~; les images
réciproques de $\R_+^*$ et $\R_-^*$ scindent $GL_n(\R)$. (Chaque
morceau est en fait \omterm{def:b2:metric:connected}{connexe} par arcs --- un exercice agréable
au-delà de nos besoins.) Par contraste $GL_n(\C)$ \emph{est} \omterm{def:b2:metric:connected}{connexe}
par arcs~: \cref{exo:b2:metric:10}.
\end{example}

\begin{example}[Un point fixe par la connexité seule]\label{ex:b2:metric:ivtfixed}
Toute application \omterm{def:b2:metric:continuity}{continue} $f \colon \intcc01 \to \intcc01$ a un
point fixe --- pas d'hypothèse de contraction, pas d'itération.
Considérer $g(x) = f(x) - x$, \omterm{def:b2:metric:continuity}{continue} sur le \omterm{def:b2:metric:connected}{connexe} $\intcc01$~:
\[
g(0) = f(0) \geq 0,
\qquad
g(1) = f(1) - 1 \leq 0 ,
\]
et le théorème des valeurs intermédiaires
(\cref{thm:b2:metric:connectedness}~(2)) livre un zéro de $g$,
c'est-à-dire un point fixe de $f$. À comparer avec Banach
(\cref{thm:b2:metric:banach})~: ici l'existence est topologique et
gratuite, mais l'unicité et l'algorithme sont perdus --- $f =
\mathrm{id}$ a tout point fixe, et l'itération d'un $f$ non
contractant peut boucler à jamais. Les deux théorèmes de point fixe
de ce chapitre répondent à des questions différentes avec des
monnaies différentes.
\end{example}

\begin{example}[$\R$ et $\R^2$ ne sont pas homéomorphes]\label{ex:b2:metric:rnotr2}
La connexité est une empreinte topologique. Supposons que $h \colon
\R^2 \to \R$ soit un homéomorphisme (une bijection \omterm{def:b2:metric:continuity}{continue} de
réciproque \omterm{def:b2:metric:continuity}{continue}). Retirer un point $a \in \R^2$~: la restriction
$h \colon \R^2\setminus\{a\} \to \R\setminus\{h(a)\}$ est encore un
homéomorphisme. Mais $\R^2$ privé d'un point est \omterm{def:b2:metric:connected}{connexe} par arcs ---
joindre deux points quelconques par un segment, en faisant un détour
par un second segment via un point auxiliaire si $a$ bloque le chemin
direct --- donc \omterm{def:b2:metric:connected}{connexe} (\cref{thm:b2:metric:connectedness}~(3))~;
tandis que $\R$ privé d'un point se scinde en deux demi-droites
\omterm{def:b2:metric:topology}{ouvertes} non vides~: non \omterm{def:b2:metric:connected}{connexe}. La connexité est préservée par les
applications \omterm{def:b2:metric:continuity}{continues}~: contradiction. Le plan et la droite sont
véritablement différents \emph{en tant qu'espaces topologiques} ---
un fait que la cardinalité seule (des bijections de type
\cref{exo:b2:structures:3} existent bien~!) est trop grossière pour
voir.
\end{example}

\section{Exercices}

\begin{exercise}[$\star$]\label{exo:b2:metric:1}
Sur $\R$, vérifier que $\delta(x, y) = \min(1, \abs{x - y})$ et
$d(x,y) = \abs{\arctan x - \arctan y}$ sont des distances. Quelles
suites convergent pour chacune~? $(\R, d)$ est-il \omterm{def:b2:metric:complete}{complet}~?
\end{exercise}

\begin{exercise}[$\star$]\label{exo:b2:metric:2}
Dans un \omterm{def:b2:metric:def}{espace métrique}, prouver qu'une suite convergente est de
Cauchy et bornée, et qu'une suite de Cauchy ayant une sous-suite
convergente converge. En déduire à nouveau que les \omterm{def:b2:metric:compact}{espaces métriques
compacts} sont \omterm{def:b2:metric:complete}{complets}.
\end{exercise}

\begin{exercise}[$\star$]\label{exo:b2:metric:3}
Dans $\bigl(C(\intcc{0}{1}), d_\infty\bigr)$, calculer la distance
entre $f(x) = x$ et $g(x) = x^2$~; décrire la boule fermée
$\overline B(0, 1)$~; et prouver que l'ensemble $\{f : f(0) = 0\}$
est fermé tandis que $\{f : f(0) > 0\}$ est \omterm{def:b2:metric:topology}{ouvert}.
\end{exercise}

\begin{exercise}[$\star\star$]\label{exo:b2:metric:4}
Prouver que l'\omterm{def:b2:metric:def}{espace métrique} discret $X$ (un ensemble quelconque)
est \omterm{def:b2:metric:complete}{complet}, et qu'il est \omterm{def:b2:metric:compact}{compact} si et seulement si $X$ est fini.
Quelles parties sont \omterm{def:b2:metric:connected}{connexes}~?
\end{exercise}

\begin{exercise}[$\star\star$]\label{exo:b2:metric:5}
Soit $X$ \omterm{def:b2:metric:compact}{compact} et $f \colon X \to X$ avec
\[
d\bigl(f(x), f(y)\bigr) < d(x, y) \quad \text{pour tous } x \neq y .
\]
Prouver que $f$ a un unique point fixe \emph{(minimiser $x \mapsto
d(x, f(x))$)}, et donner un exemple sur $X = \intco{1}{+\infty}$
(non \omterm{def:b2:metric:compact}{compact}) sans point fixe.
\end{exercise}

\begin{exercise}[$\star\star$]\label{exo:b2:metric:6}
(\omterm{def:b2:metric:compact}{Compacts} emboîtés) Soit $(K_n)$ une suite décroissante de parties
\omterm{def:b2:metric:compact}{compactes} non vides d'un \omterm{def:b2:metric:def}{espace métrique}. Prouver que $\bigcap_n K_n
\neq \emptyset$ \emph{(choisir $x_n \in K_n$ et extraire)}. Montrer
par un exemple que des \emph{fermés} emboîtés non vides dans $\R$
peuvent avoir une intersection vide.
\end{exercise}

\begin{exercise}[$\star\star$]\label{exo:b2:metric:7}
Soit $K$ \omterm{def:b2:metric:compact}{compact} et $f \colon K \to Y$ \omterm{def:b2:metric:continuity}{continue} et bijective.
Prouver que $f^{-1}$ est \omterm{def:b2:metric:continuity}{continue} \emph{(utiliser les fermés~:
\cref{thm:b2:metric:globalcontinuity} et
\cref{thm:b2:metric:compactprops})}. Donner un contre-exemple sans
compacité ($\gamma(t) = (\cos t, \sin t)$ sur $\intco{0}{2\pi}$).
\end{exercise}

\begin{exercise}[$\star\star$]\label{exo:b2:metric:8}
L'\emph{ensemble de Cantor} $C$ s'obtient à partir de $\intcc{0}{1}$
en supprimant répétitivement les tiers centraux \omterm{def:b2:metric:topology}{ouverts}. Prouver que
$C$ est \omterm{def:b2:metric:compact}{compact}, d'intérieur vide, et infini --- en fait équipotent à
$\{0,1\}^{\N}$ \emph{(développements en base trois avec chiffres
$0,2$~; \cref{exo:b2:structures:3})}.
\end{exercise}

\begin{exercise}[$\star\star\star$]\label{exo:b2:metric:9}
Soit $X$ un \omterm{def:b2:metric:def}{espace métrique} \emph{\omterm{def:b2:metric:compact}{compact}} et $f \colon X \to X$ une
isométrie~: $d(f(x), f(y)) = d(x, y)$. Prouver que $f$ est
surjective. \emph{Indication~: si $a \notin f(X)$, alors $\varepsilon
= d(a, f(X)) > 0$ (pourquoi~?)~; étudier l'orbite $a, f(a), f^2(a),
\dots$ et montrer que ses points sont deux à deux distants de $\geq
\varepsilon$ --- contradiction avec la compacité.}
\end{exercise}

\begin{exercise}[$\star\star\star$]\label{exo:b2:metric:10}
Prouver que $GL_n(\C)$ est \omterm{def:b2:metric:connected}{connexe} par arcs. \emph{Indication~:
étant donné $A, B$ inversibles, considérer $p(z) = \det\bigl((1 -
z)A + zB\bigr)$ pour $z \in \C$~: un polynôme en $z$, non
identiquement nul, donc à racines en nombre fini~; joindre $0$ à $1$
dans $\C$ par un chemin les évitant.}
\end{exercise}

\begin{exercise}[$\star\star$]\label{exo:b2:metric:11}
Pour $\emptyset \neq A \subseteq X$, poser $d(x, A) = \inf_{a \in A}
d(x, a)$. Prouver que $x \mapsto d(x, A)$ est $1$-lipschitzienne, que
$d(x, A) = 0$ si et seulement si $x \in \overline A$, et que pour des
\emph{fermés} disjoints non vides $A, B$ la fonction
\[
\varphi(x) = \frac{d(x, A)}{d(x, A) + d(x, B)}
\]
est bien définie, \omterm{def:b2:metric:continuity}{continue}, égale à $0$ exactement sur $A$ et à $1$
exactement sur $B$ --- un «~interrupteur~» \omterm{def:b2:metric:continuity}{continu} séparant deux
fermés disjoints quelconques.
\end{exercise}

\begin{exercise}[$\star\star\star$]\label{exo:b2:metric:12}
(Baire) Soit $X$ un \omterm{def:b2:metric:complete}{espace métrique complet} et $(U_n)_{n\geq1}$ une
suite de parties \omterm{def:b2:metric:topology}{ouvertes} denses. Prouver que $\bigcap_n U_n$ est
dense dans $X$ \emph{(à l'intérieur d'une boule quelconque,
construire des boules fermées emboîtées $\overline B(x_n, r_n)
\subseteq U_n$ avec $r_n \to 0$ et utiliser la complétude)}. En
déduire que $\R$ n'est pas une réunion \omterm{def:b2:structures:countable}{dénombrable} de fermés
d'intérieur vide, et retrouver --- encore --- que $\R$ est non
\omterm{def:b2:structures:countable}{dénombrable}.
\end{exercise}

\section{Problème~: itération de Picard}

Complétude plus contraction est une machine à résoudre~: donnez-lui
une équation écrite comme un problème de point fixe, et elle renvoie
existence, unicité, un algorithme, et des barres d'erreur. Ce problème
de fin de semaine fait tourner la machine à pleine puissance sur
l'équation $y' = f(t, y)$~: nous prouvons le \emph{théorème local de
Picard--Lindelöf} (le cœur non linéaire de la théorie de
Cauchy--Lipschitz du~\cref{ch:b2:diffeq}), observons chaque hypothèse
gagner son pain à travers des contre-exemples, et récoltons des
dividendes purement métriques --- dépendance \omterm{def:b2:metric:continuity}{continue} par rapport aux
données, équation de Kepler, et auto-similarité de l'ensemble de
Cantor.

\begin{omfigure}
\begin{tikzpicture}
\begin{axis}[omaxis, width=8.6cm, height=5.4cm,
    xmin=-0.05, xmax=1.45, ymin=0, ymax=4.2,
    xtick={0,1}, ytick={1,2,3}]
  \addplot[omDef!45, thick, domain=0:1.4, samples=2] {1};
  \addplot[omDef!65, thick, domain=0:1.4, samples=2] {1 + x};
  \addplot[omDef, thick, domain=0:1.4, samples=60]
    {1 + x + x^2/2};
  \addplot[omThm, very thick, domain=0:1.4, samples=120]
    {exp(x)};
  \node[font=\small, omDef] at (axis cs:1.28,0.8) {$y_0$};
  \node[font=\small, omDef] at (axis cs:1.28,2.15) {$y_1$};
  \node[font=\small, omDef] at (axis cs:1.28,2.95) {$y_2$};
  \node[font=\small, omThm] at (axis cs:1.17,3.85)
    {$\eu^{\,t}$};
\end{axis}
\end{tikzpicture}

{\small Itérés de Picard pour $y' = y$, $y(0) = 1$~: chaque passage
par $T(y)(t) = 1 + \int_0^t y$ ajoute un terme de Taylor, et la
contraction resserre toute la suite uniformément sur $\eu^{\,t}$.}
\end{omfigure}

\begin{problem}[{Problème de fin de semaine --- le théorème de
Picard--Lindelöf}]
\label{pb:b2:metric:1}
Tout au long, $t_0 \in \R$, $y_0 \in \R$, $a, b > 0$, et $f$ est une
fonction \omterm{def:b2:metric:continuity}{continue} sur le rectangle $R = \intcc{t_0 - a}{t_0 + a}
\times \intcc{y_0 - b}{y_0 + b}$, bornée par $M = \sup_R\,\abs f$, et
\emph{$L$-lipschitzienne en sa seconde variable}~: $\abs{f(t, y) -
f(t, z)} \leq L\abs{y - z}$ dès que les deux points sont dans $R$.
Poser
\[
h = \min\Bigl(a, \frac bM\Bigr) \quad (\text{avec } h = a
\text{ si } M = 0), \qquad I = \intcc{t_0 - h}{t_0 + h}.
\]

\medskip
\noindent\textbf{Partie I --- L'étage \omterm{def:b2:metric:complete}{complet}.}
\begin{enumerate}
  \item Prouver les deux assertions citées dans la~%
        \cref{def:b2:metric:complete}~: une partie fermée d'un \omterm{def:b2:metric:complete}{espace
        métrique complet} est complète, et une partie complète d'un
        \omterm{def:b2:metric:def}{espace métrique} quelconque est fermée. En déduire que toute
        partie fermée de $\bigl(C(I), d_\infty\bigr)$ est un \omterm{def:b2:metric:complete}{espace
        métrique complet}.
  \item Montrer que l'application $C(I) \to C(I)$, $y \mapsto \bigl(t
        \mapsto \int_{t_0}^{t}y(s)\,\dd s\bigr)$, est
        $h$-lipschitzienne pour $d_\infty$.
  \item (Les points fixes bougent moins que les applications) Soit $g
        \colon X \to X$ une $k$-contraction d'un \omterm{def:b2:metric:def}{espace métrique} de
        point fixe $\ell_g$, et $\widetilde g \colon X \to X$ une
        application \emph{quelconque} ayant un point fixe
        $\ell_{\widetilde g}$. Prouver
        \[
        d(\ell_g, \ell_{\widetilde g}) \leq
        \frac{d\bigl(g(\ell_{\widetilde g}),
        \widetilde g(\ell_{\widetilde g})\bigr)}{1 - k}
        \leq \frac{\sup_{x \in X} d\bigl(g(x), \widetilde
        g(x)\bigr)}{1 - k}.
        \]
  \item (L'astuce de l'itéré) Soit $X$ \omterm{def:b2:metric:complete}{complet} non vide et $g \colon
        X \to X$ une application --- non supposée \omterm{def:b2:metric:continuity}{continue} --- telle
        qu'un certain itéré $g^m$ soit une $k$-contraction. Prouver
        que $g$ a un unique point fixe $\ell$ et que \emph{toute}
        orbite $x_{n+1} = g(x_n)$ converge vers $\ell$. \emph{(Les
        points fixes de $g$ sont des points fixes de $g^m$~;
        réciproquement $g(\ell)$ est un point fixe de $g^m$~;
        découper l'orbite selon les résidus mod $m$.)}
\end{enumerate}

\medskip
\noindent\textbf{Partie II --- Le théorème de Picard--Lindelöf.}
\begin{enumerate}[resume]
  \item Montrer qu'une fonction $y \colon I \to \intcc{y_0 - b}{y_0 +
        b}$ est $C^1$ avec $y(t_0) = y_0$ et $y' = f(t, y)$ sur $I$
        si et seulement si elle est \omterm{def:b2:metric:continuity}{continue} et satisfait l'équation
        intégrale
        \[
        y(t) = y_0 + \int_{t_0}^{t} f\bigl(s, y(s)\bigr)\dd s
        \qquad (t \in I).
        \]
  \item Soit $X_h = \{y \in C(I) : \abs{y(t) - y_0} \leq b
        \text{ sur } I\}$ et soit $T$ défini par $T(y)(t) = y_0 +
        \int_{t_0}^{t}f(s, y(s))\dd s$. Montrer que $X_h$ est une
        partie fermée non vide de $C(I)$, donc complète, et que $T$
        envoie $X_h$ dans $X_h$ --- c'est ici que $h \leq b/M$
        opère.
  \item Montrer que $d_\infty\bigl(T(y), T(z)\bigr) \leq
        Lh\,d_\infty(y, z)$ sur $X_h$~: si $Lh < 1$, le théorème de
        Banach conclut déjà. Nous supprimons cette condition de
        petitesse ensuite.
  \item Prouver par récurrence sur $n$~:
        \[
        \abs{T^n(y)(t) - T^n(z)(t)} \leq
        \frac{\bigl(L\abs{t - t_0}\bigr)^n}{n!}\,
        d_\infty(y, z) \qquad (y, z \in X_h,\ t \in I),
        \]
        de sorte qu'un certain itéré de $T$ est une contraction.
        Conclure avec la question 4 (\emph{le théorème de
        Picard--Lindelöf})~: le problème de Cauchy $y' = f(t,y)$,
        $y(t_0) = y_0$ a exactement une solution sur $I =
        \intcc{t_0 - h}{t_0 + h}$ à valeurs dans $\intcc{y_0 -
        b}{y_0 + b}$.
  \item Montrer que la restriction «~à valeurs dans $\intcc{y_0 -
        b}{y_0 + b}$~» est automatique~: toute solution du problème
        de Cauchy définie sur $I$ dont le graphe part dans $R$ reste
        dans $\intcc{y_0 - b}{y_0 + b}$ \emph{(considérer le premier
        instant de sortie et majorer $\abs{y(t) - y_0}$ par
        $M\abs{t - t_0}$)}. Donc l'unicité vaut parmi toutes les
        solutions sur $I$.
  \item Faire tourner la machine sur $y' = y$, $y(0) = 1$, en
        partant de la constante $y^{(0)} \equiv 1$~: calculer les
        itérés de Picard $y^{(n)}$, les identifier, et décrire la
        convergence.
\end{enumerate}

\medskip
\noindent\textbf{Partie III --- Chaque hypothèse gagne son pain.}
\begin{enumerate}[resume]
  \item (Lipschitz échoue, l'unicité échoue) Pour $y' =
        2\sqrt{\abs y}$, $y(0) = 0$~: vérifier que $y \equiv 0$ et,
        pour tout $c \geq 0$, la fonction $y_c(t) = 0$ pour $t \leq
        c$, $y_c(t) = (t - c)^2$ pour $t > c$, sont toutes des
        solutions $C^1$ sur $\R$. Où exactement $y \mapsto
        2\sqrt{\abs y}$ cesse-t-il d'être \omterm{def:b2:metric:continuity}{lipschitzien}~?
  \item (La localité est réelle) Pour $y' = y^2$, $y(0) = 1$~:
        résoudre explicitement, donner l'intervalle maximal
        d'existence, et calculer le meilleur $h$ que le théorème
        peut certifier sur tous les choix du rectangle ($a$ grand,
        $b$ libre)~: montrer $h_{\max} = \sup_{b>0}
        \frac{b}{(1+b)^2} = \frac14$, tandis que la vraie solution
        vit sur $\intoo{-\infty}{1}$.
  \item (La complétude n'est pas du décor) Sur $X = \Q \cap
        \intcc{1}{2}$ avec la distance usuelle, soit $g(x) = \frac
        x2 + \frac1x$. Montrer que $g(X) \subseteq X$, que $g$ est
        une $\frac12$-contraction \emph{(inégalité des accroissements
        finis)}, et que $g$ n'a aucun point fixe dans $X$. Quelle
        hypothèse du théorème de Banach échoue, et quel est le point
        fixe dans le complété~?
  \item (Barres d'erreur) Pour une $k$-contraction $g$ sur un espace
        \omterm{def:b2:metric:complete}{complet}, prouver l'estimation \emph{a posteriori} $d(x_n,
        \ell) \leq \frac{k}{1-k}\,d(x_n, x_{n-1})$. Pour
        l'application de Héron $g(x) = \frac x2 + \frac 1x$ sur
        $\intcc{1}{2}$ (point fixe $\sqrt2$), en partant de $x_0 =
        \frac32$~: combien d'étapes la majoration \emph{a priori}
        $\frac{k^n}{1-k}d(x_1, x_0)$ exige-t-elle pour une précision
        $10^{-6}$, et combien d'étapes suffisent en réalité~?
        (Calculer $x_1, x_2, x_3$ et leurs erreurs~; la majoration
        par contraction est honnête mais pessimiste --- Héron
        converge quadratiquement.)
\end{enumerate}

\medskip
\noindent\textbf{Partie IV --- Dépendance \omterm{def:b2:metric:continuity}{continue}.} Dans cette
partie $Lh < 1$, donc $T$ lui-même est une contraction sur $X_h$
(question 7)~; on note $y[\,y_0\,]$ la solution de valeur initiale
$y_0$.
\begin{enumerate}[resume]
  \item (Dépendance en la valeur initiale) Soit $z_0$ une autre
        valeur initiale avec $\abs{z_0 - y_0}$ assez petit pour que
        les deux problèmes tiennent dans le rectangle. En utilisant
        la question 3, prouver
        \[
        d_\infty\bigl(y[y_0], y[z_0]\bigr) \leq
        \frac{\abs{y_0 - z_0}}{1 - Lh}.
        \]
  \item (Longs intervalles par chaînage) Supposons que les solutions
        existent sur un long segment découpé en $m$ morceaux
        consécutifs sur chacun desquels la majoration précédente
        s'applique avec $Lh \leq \frac12$. Montrer que l'écart croît
        d'un facteur au plus $2$ par morceau, donc $d_\infty \leq
        2^m\abs{y_0 - z_0}$ au total --- une majoration exponentielle
        en la longueur, l'ombre discrète du $\eu^{L\abs{t - t_0}}$
        du lemme de Gronwall (\cref{ch:b2:diffeq}).
  \item (Dépendance en le champ) Soit $g$ un autre champ sur $R$,
        également $L$-lipschitzien en $y$, avec $\sup_R \abs{f - g}
        \leq \varepsilon$. Prouver que les solutions correspondantes
        satisfont $d_\infty \leq \frac{\varepsilon h}{1 - Lh}$~:
        l'erreur de modélisation se propage linéairement.
  \item (Paramètres) Si une famille $f_\lambda$ de champs est
        $L$-lipschitzienne en $y$ uniformément et $\sup_R\abs{f_\lambda
        - f_\mu} \leq C\abs{\lambda - \mu}$, en déduire que $\lambda
        \mapsto y_\lambda$ est \omterm{def:b2:metric:continuity}{lipschitzienne} de l'espace des
        paramètres vers $\bigl(C(I), d_\infty\bigr)$.
  \item (Les systèmes ne coûtent rien) Expliquer pourquoi les
        parties I, II et IV valent mot pour mot pour $y$ à valeurs
        dans $\R^n$ (distances sup bâties sur l'une des distances du~%
        \cref{ex:b2:metric:examples}), puis calculer tous les itérés
        de Picard pour le système $y' = Ay$, $y(0) = (c_1, c_2)$, $A
        = \begin{psmallmatrix} 0 & 1\\ 0 & 0\end{psmallmatrix}$~:
        montrer que l'itération devient stationnaire à la solution
        exacte après une étape.
\end{enumerate}

\medskip
\noindent\textbf{Partie V --- Dividendes métriques et synthèse.}
\begin{enumerate}[resume]
  \item (Perturbation de l'identité) Soit $X$ un étage \omterm{def:b2:metric:complete}{complet} de
        type espace normé~: prendre $X = C(I)$ ou $\R^n$. Si $\eta
        \colon X \to X$ est $k$-lipschitzienne avec $k < 1$, prouver
        que $x \mapsto x + \eta(x)$ est une bijection de $X$ dont la
        réciproque est $\frac1{1-k}$-lipschitzienne \emph{(pour
        chaque $y$, appliquer Banach à $x \mapsto y - \eta(x)$)}.
        C'est le cœur métrique du théorème d'inversion locale
        (\cref{ch:b2:diffcalc}).
  \item (Équation de Kepler) Pour $0 \leq e < 1$ et $m \in \R$,
        prouver que $x = m + e\sin x$ a exactement une solution, que
        l'itération $x_{n+1} = m + e\sin x_n$ converge vers elle
        depuis tout point de départ, et estimer~: pour $e =
        \frac12$, $m = 1$, combien d'itérations garantissent une
        erreur $\leq 10^{-3}$ par la majoration a priori~? (La
        solution est $x \approx 1.4987$.)
  \item (L'ensemble de Cantor est un point fixe) Soit $S_1(x) =
        \frac x3$ et $S_2(x) = \frac x3 + \frac23$ sur $\R$, et soit
        $C$ l'ensemble de Cantor du~\cref{exo:b2:metric:8}. Prouver
        $C = S_1(C) \cup S_2(C)$, et expliquer en une phrase
        pourquoi aucun \emph{autre} ensemble \omterm{def:b2:metric:compact}{compact} non vide ne
        satisfait cette équation (l'application $A \mapsto S_1(A)
        \cup S_2(A)$ est une contraction pour une distance entre
        \omterm{def:b2:metric:compact}{compacts} --- la distance de Hausdorff, rendue rigoureuse
        dans le volume de troisième année).
  \item (La connexité globalise l'unicité) Soit $f$ localement
        \omterm{def:b2:metric:continuity}{lipschitzienne} en $y$ sur un \omterm{def:b2:metric:topology}{ouvert}, et soient $y, z$ deux
        solutions de $y' = f(t, y)$ sur un intervalle commun $J$
        avec $y(t_0) = z(t_0)$. Prouver $y = z$ sur $J$~: montrer que
        $\{t \in J : y(t) = z(t)\}$ est non vide, fermé dans $J$, et
        \omterm{def:b2:metric:topology}{ouvert} dans $J$ (par unicité locale), et utiliser la
        connexité des intervalles
        (\cref{thm:b2:metric:connectedness}).
  \item (Pas de petitesse pour les équations linéaires) Pour $y' =
        \alpha(t)y + \beta(t)$ avec $\alpha, \beta$ \omterm{def:b2:metric:continuity}{continues} sur un
        segment $\intcc{A}{B}$, adapter la question 8 pour montrer
        que la majoration factorielle vaut sur le segment
        \emph{entier}, donc existence et unicité sont globales
        là-dessus --- le cas scalaire du théorème de
        Cauchy--Lipschitz du~\cref{ch:b2:diffeq}, sans restriction
        sur la longueur $B - A$.
  \item (Synthèse) Une phrase pour chacun~: ce qu'a apporté la
        complétude~; ce qu'a apporté la contraction~; ce qu'a
        acheté l'astuce de l'itéré par rapport au Banach ordinaire~;
        où la connexité est entrée~; et quel contre-exemple de la
        partie III garde quelle hypothèse. Nommer le théorème sommet,
        et dire ce qui remplace la contraction lorsque $f$ est
        seulement \omterm{def:b2:metric:continuity}{continue} (le théorème de Peano, via la compacité
        dans les espaces de fonctions --- l'Arzelà--Ascoli du volume
        de troisième année).
\end{enumerate}
\end{problem}
